\documentclass[12pt]{article}
\usepackage[utf8]{inputenc}
\usepackage[T1]{fontenc}
\usepackage{geometry}
\geometry{margin=2.5cm}
\usepackage{enumitem}
\usepackage{setspace}
\onehalfspacing

\providecommand{\AlunoNome}{\textit{Aluno não informado}}

\title{Prova Semestral de Testes de Software}
\author{Centro Universitário Senac \\
Professor: Afonso Brandão}
\date{Disciplina: Testes de Software}

\begin{document}

\maketitle

\begin{center}
  \textbf{Aluno(a):} \AlunoNome
\end{center}

\begin{center}
  \textbf{Instruções:}
\end{center}
\begin{itemize}[leftmargin=*]
  \item Leia com atenção e responda de forma clara e organizada.
  \item A resposta deve ser redigida em texto corrido, destacando cada tópico estudado.
  \item Utilize exemplos ou referências às práticas realizadas ao longo do semestre.
\end{itemize}

\section*{Questão Única}
\textbf{Contexto:} o objetivo desta avaliação é permitir que você reflita sobre o conteúdo visto durante o semestre e demonstre a compreensão que construiu.

\textbf{Enunciado:} elabore um texto dissertativo no qual você apresente o que compreendeu ao longo do semestre de Testes de Software. \textit{Mostre todos os tópicos estudados}, agrupando-os em blocos temáticos ou por ordem cronológica das aulas, indicando a relevância de cada um para sua formação. Cite técnicas, ferramentas, tipos de teste, estratégias de automação, métricas e qualquer outra noção abordada em sala ou nas atividades práticas. Finalize comentando sobre como pode aplicar esse repertório em situações reais de desenvolvimento.

\textbf{Critérios de avaliação:}
\begin{enumerate}[label=\arabic*.]
  \item Compreensão dos tópicos e sua articulação em um texto coerente.
  \item Capacidade de citar exemplos e conectar teoria e prática.
  \item Organização e clareza da argumentação, incluindo a conclusão sobre aplicação futura.
\end{enumerate}

\end{document}
